% ============================================================================
% Fichier  : annexes/Ann_G_exigences.tex
% Version  : 1.0 -- novembre 2026
% Objet    : les quarante-six exigences de la mission NYANGA, par domaine,
%            avec leur cote, leur preuve et la nature de chaque NEK.
%            Les lignes reprennent mot pour mot les encadrés des chapitres
%            8, 11, 12, 13, 15 et 16.
% ============================================================================

% Évite un titre de section isolé en bas de page.
\providecommand{\gardetitre}{\par\vskip 0pt plus 9\baselineskip\penalty-300\vskip 0pt plus -9\baselineskip}

\chapter{Les Quarante-six Exigences de la Mission NYANGA}
\label{ann:exigences}

Cette annexe réunit les exigences cotées par KORO AUDIT à NYANGA CRÉDIT du
24~juin au 25~septembre~2026. Elle permet de refaire chaque calcul du
livre~: 8 + 10 + 8 + 8 + 8 + 4 = 46 exigences, 8~\OUI{}, 28~\NON{},
10~\NEK{}. La dernière colonne donne la nature de chaque \NEK{}~: de
\textit{preuve} quand une pièce manque, de \textit{critère} quand c'est
le texte ou la référence qui manque.

\gardetitre
\section*{D1. Recensement et fondements (Chapitre~\ref{chap:audit})}

\begin{tableaucours}{cL{6.2cm}cL{3.6cm}}{%
    \tableauHeader{N°} &
    \tableauHeader{Exigence} &
    \tableauHeader{Cote} &
    \tableauHeader{Preuve, nature du NEK}
}
1 & Les familles de données sont recensées & \OUI{} & Inventaire du 7~juillet \\
2 & Un registre des activités de traitement est tenu & \NON{} & Aucun registre \\
3 & Chaque traitement a un fondement valable & \NON{} & Formulaires, cahiers \\
4 & L'autorisation de l'Autorité est demandée & \NON{} & Attestation \\
5 & Une durée de conservation est fixée & \NON{} & Aucune règle \\
6 & L'information de l'article~21 est fournie & \NON{} & Formulaire 2016 \\
7 & L'identification des clients repose sur une obligation légale & \NEK{} &
Texte non produit. \textit{Critère} \\
8 & La vidéosurveillance repose sur un fondement valable & \NEK{} &
Fondement incertain. \textit{Critère} \\
\end{tableaucours}

\gardetitre
\section*{D2. Droits des personnes (Chapitre~\ref{chap:droits})}

\begin{tableaucours}{cL{6.2cm}cL{3.6cm}}{%
    \tableauHeader{N°} &
    \tableauHeader{Exigence} &
    \tableauHeader{Cote} &
    \tableauHeader{Preuve, nature du NEK}
}
9 & Les personnes sont informées de leurs droits & \NON{} & Formulaire 2016 \\
10 & Une procédure de traitement des demandes existe & \NON{} & Aucune \\
11 & Un registre des demandes est tenu & \NON{} & Reconstitué par l'audit \\
12 & Les demandes d'accès reçoivent une réponse & \NON{} & 12 sur 31 \\
13 & Les réponses respectent le critère de 30 jours & \NON{} & 8 sur 12 \\
14 & Les oppositions sont exécutées dans les systèmes & \NON{} & 3 numéros
encore appelés \\
15 & La portabilité est possible & \NON{} & Refus du 14~août \\
16 & L'identité est vérifiée de façon proportionnée & \NEK{} & Pratiques non
documentées. \textit{Preuve} \\
17 & Les effacements acceptés sont exécutés & \NEK{} & Sauvegardes non
examinées. \textit{Preuve} \\
18 & La procédure de succession respecte l'article~45 & \OUI{} & Note de
2025, 3 dossiers \\
\end{tableaucours}

\gardetitre
\section*{D3. Sécurité et confidentialité (Chapitre~\ref{chap:securite})}

\begin{tableaucours}{cL{6.2cm}cL{3.6cm}}{%
    \tableauHeader{N°} &
    \tableauHeader{Exigence} &
    \tableauHeader{Cote} &
    \tableauHeader{Preuve, nature du NEK}
}
19 & L'équipement à l'origine de l'incident est remplacé et tenu à jour &
\OUI{} & Bon de livraison, état des correctifs \\
20 & Un inventaire des équipements est tenu & \OUI{} & Inventaire de
mars~2026 \\
21 & Les journaux sont centralisés et conservés & \OUI{} & Console, treize
mois en ligne \\
22 & Les sauvegardes sont restaurables & \OUI{} & Test du 3~septembre \\
23 & Chaque agent n'accède qu'aux clients de son agence & \NON{} & Test sur
six comptes \\
24 & Les relevés sont transmis par un canal maîtrisé & \NON{} & Messagerie
personnelle \\
25 & Les habilitations sont revues & \NON{} & Aucune revue depuis 2023 \\
26 & Les mesures suivent le référentiel de l'Autorité (art.~24) & \NEK{} &
Référentiel non publié. \textit{Critère} \\
\end{tableaucours}

\gardetitre
\section*{D4. Sous-traitants et transferts (Chapitre~\ref{chap:soustraitants})}

\begin{tableaucours}{cL{6.2cm}cL{3.6cm}}{%
    \tableauHeader{N°} &
    \tableauHeader{Exigence} &
    \tableauHeader{Cote} &
    \tableauHeader{Preuve, nature du NEK}
}
27 & Chaque sous-traitant est lié par un contrat & \OUI{} & Trois contrats \\
28 & Les contrats définissent données et personnes (art.~30(3)) & \NON{} &
SOREC, opérateur \\
29 & Les contrats contiennent les douze clauses & \NON{} & Un sur trois au
mieux \\
30 & Le droit d'audit est prévu & \NON{} & Aucun contrat \\
31 & SOREC détruit les données des dossiers clos & \NON{} & Classeurs
conservés \\
32 & L'éditeur applique les mesures de sécurité convenues & \NEK{} & Ni
audit ni rapport. \textit{Preuve} \\
33 & L'opérateur de SMS limite l'usage des numéros & \NEK{} & Aucune pièce.
\textit{Preuve} \\
34 & La convention avec SAWA ASSURANCES est publique & \NON{} & Inexistante \\
\end{tableaucours}

\gardetitre
\section*{D5. Violations (Chapitre~\ref{chap:violations})}

\begin{tableaucours}{cL{6.2cm}cL{3.6cm}}{%
    \tableauHeader{N°} &
    \tableauHeader{Exigence} &
    \tableauHeader{Cote} &
    \tableauHeader{Preuve, nature du NEK}
}
35 & Une procédure de gestion des violations existe & \NON{} & Aucune \\
36 & Un registre des violations est tenu & \NON{} & Aucun \\
37 & La violation de 2025 est qualifiée pour les personnes & \NON{} &
Rapport technique seul \\
38 & La demande de différer est écrite et motivée & \OUI{} & Lettre du
20~mars~2025 \\
39 & La demande a été réexaminée & \NON{} & Aucune démarche \\
40 & Les personnes ont été informées quand la raison a cessé & \NON{} &
Rien au 31~août~2026 \\
41 & L'Autorité a été informée & \NEK{} & Autorité non organisée.
\textit{Critère} \\
42 & L'obligation d'informer s'appliquait en mars~2025 & \NEK{} & Question
de droit non tranchée. \textit{Critère} \\
\end{tableaucours}

\gardetitre
\section*{D6. Gouvernance (Chapitre~\ref{chap:programme})}

\begin{tableaucours}{cL{6.2cm}cL{3.6cm}}{%
    \tableauHeader{N°} &
    \tableauHeader{Exigence} &
    \tableauHeader{Cote} &
    \tableauHeader{Preuve, nature du NEK}
}
43 & Le risque de chaque traitement est évalué (art.~25) & \NON{} & Aucune
évaluation \\
44 & Une évaluation d'impact précède les traitements à risque élevé
(art.~33) & \NON{} & Douze traitements, aucune \\
45 & Un dossier de conformité peut être présenté au contrôle (art.~35) &
\NON{} & Aucun dossier \\
46 & Le rapport annuel sur les mesures est adressé à l'Autorité
(art.~27(2)) & \NEK{} & Autorité non organisée. \textit{Critère} \\
\end{tableaucours}

\gardetitre
\section*{Récapitulation}

\begin{tableaucours}{cL{4.2cm}ccccc}{%
    \tableauHeader{Code} &
    \tableauHeader{Domaine} &
    \tableauHeader{\OUI{}} &
    \tableauHeader{\NON{}} &
    \tableauHeader{\NEK{}} &
    \tableauHeader{de preuve} &
    \tableauHeader{de critère}
}
D1 & Recensement et fondements & 1 & 5 & 2 & 0 & 2 \\
D2 & Droits des personnes & 1 & 7 & 2 & 2 & 0 \\
D3 & Sécurité et confidentialité & 4 & 3 & 1 & 0 & 1 \\
D4 & Sous-traitants et transferts & 1 & 5 & 2 & 2 & 0 \\
D5 & Violations & 1 & 5 & 2 & 0 & 2 \\
D6 & Gouvernance & 0 & 3 & 1 & 0 & 1 \\
Total & & 8 & 28 & 10 & 4 & 6 \\
\end{tableaucours}

Les quatre \NEK{} de preuve se lèvent par le travail de l'auditeur ou de
l'établissement~: une procédure écrite, l'examen des sauvegardes, un audit
des deux prestataires. Les six \NEK{} de critère ne dépendent ni de l'un ni
de l'autre. Ils attendent un texte d'application, l'organisation de
l'Autorité ou une décision de justice
(Chapitre~\ref{chap:recherche}). Le transfert vers Abidjan n'entre pas dans
ce décompte~: le Chapitre~\ref{chap:programme} dit pourquoi et le cote à
part.
